% Options for packages loaded elsewhere
% Options for packages loaded elsewhere
\PassOptionsToPackage{unicode}{hyperref}
\PassOptionsToPackage{hyphens}{url}
%
\documentclass[
  11pt,
  ignorenonframetext,
  aspectratio=169,
]{beamer}
\newif\ifbibliography
\usepackage{pgfpages}
\setbeamertemplate{caption}[numbered]
\setbeamertemplate{caption label separator}{: }
\setbeamercolor{caption name}{fg=normal text.fg}
\beamertemplatenavigationsymbolsempty
% remove section numbering
\setbeamertemplate{part page}{
  \centering
  \begin{beamercolorbox}[sep=16pt,center]{part title}
    \usebeamerfont{part title}\insertpart\par
  \end{beamercolorbox}
}
\setbeamertemplate{section page}{
  \centering
  \begin{beamercolorbox}[sep=12pt,center]{section title}
    \usebeamerfont{section title}\insertsection\par
  \end{beamercolorbox}
}
\setbeamertemplate{subsection page}{
  \centering
  \begin{beamercolorbox}[sep=8pt,center]{subsection title}
    \usebeamerfont{subsection title}\insertsubsection\par
  \end{beamercolorbox}
}
% Prevent slide breaks in the middle of a paragraph
\widowpenalties 1 10000
\raggedbottom
\AtBeginPart{
  \frame{\partpage}
}
\AtBeginSection{
  \ifbibliography
  \else
    \frame{\sectionpage}
  \fi
}
\AtBeginSubsection{
  \frame{\subsectionpage}
}
\usepackage{iftex}
\ifPDFTeX
  \usepackage[T1]{fontenc}
  \usepackage[utf8]{inputenc}
  \usepackage{textcomp} % provide euro and other symbols
\else % if luatex or xetex
  \usepackage{unicode-math} % this also loads fontspec
  \defaultfontfeatures{Scale=MatchLowercase}
  \defaultfontfeatures[\rmfamily]{Ligatures=TeX,Scale=1}
\fi
\usepackage{lmodern}

\usetheme[]{Madrid}
\usecolortheme[]{default}
\ifPDFTeX\else
  % xetex/luatex font selection
\fi
% Use upquote if available, for straight quotes in verbatim environments
\IfFileExists{upquote.sty}{\usepackage{upquote}}{}
\IfFileExists{microtype.sty}{% use microtype if available
  \usepackage[]{microtype}
  \UseMicrotypeSet[protrusion]{basicmath} % disable protrusion for tt fonts
}{}
\makeatletter
\@ifundefined{KOMAClassName}{% if non-KOMA class
  \IfFileExists{parskip.sty}{%
    \usepackage{parskip}
  }{% else
    \setlength{\parindent}{0pt}
    \setlength{\parskip}{6pt plus 2pt minus 1pt}}
}{% if KOMA class
  \KOMAoptions{parskip=half}}
\makeatother

\usepackage{color}
\usepackage{fancyvrb}
\newcommand{\VerbBar}{|}
\newcommand{\VERB}{\Verb[commandchars=\\\{\}]}
\DefineVerbatimEnvironment{Highlighting}{Verbatim}{commandchars=\\\{\}}
% Add ',fontsize=\small' for more characters per line
\usepackage{framed}
\definecolor{shadecolor}{RGB}{241,243,245}
\newenvironment{Shaded}{\begin{snugshade}}{\end{snugshade}}
\newcommand{\AlertTok}[1]{\textcolor[rgb]{0.68,0.00,0.00}{#1}}
\newcommand{\AnnotationTok}[1]{\textcolor[rgb]{0.37,0.37,0.37}{#1}}
\newcommand{\AttributeTok}[1]{\textcolor[rgb]{0.40,0.45,0.13}{#1}}
\newcommand{\BaseNTok}[1]{\textcolor[rgb]{0.68,0.00,0.00}{#1}}
\newcommand{\BuiltInTok}[1]{\textcolor[rgb]{0.00,0.23,0.31}{#1}}
\newcommand{\CharTok}[1]{\textcolor[rgb]{0.13,0.47,0.30}{#1}}
\newcommand{\CommentTok}[1]{\textcolor[rgb]{0.37,0.37,0.37}{#1}}
\newcommand{\CommentVarTok}[1]{\textcolor[rgb]{0.37,0.37,0.37}{\textit{#1}}}
\newcommand{\ConstantTok}[1]{\textcolor[rgb]{0.56,0.35,0.01}{#1}}
\newcommand{\ControlFlowTok}[1]{\textcolor[rgb]{0.00,0.23,0.31}{\textbf{#1}}}
\newcommand{\DataTypeTok}[1]{\textcolor[rgb]{0.68,0.00,0.00}{#1}}
\newcommand{\DecValTok}[1]{\textcolor[rgb]{0.68,0.00,0.00}{#1}}
\newcommand{\DocumentationTok}[1]{\textcolor[rgb]{0.37,0.37,0.37}{\textit{#1}}}
\newcommand{\ErrorTok}[1]{\textcolor[rgb]{0.68,0.00,0.00}{#1}}
\newcommand{\ExtensionTok}[1]{\textcolor[rgb]{0.00,0.23,0.31}{#1}}
\newcommand{\FloatTok}[1]{\textcolor[rgb]{0.68,0.00,0.00}{#1}}
\newcommand{\FunctionTok}[1]{\textcolor[rgb]{0.28,0.35,0.67}{#1}}
\newcommand{\ImportTok}[1]{\textcolor[rgb]{0.00,0.46,0.62}{#1}}
\newcommand{\InformationTok}[1]{\textcolor[rgb]{0.37,0.37,0.37}{#1}}
\newcommand{\KeywordTok}[1]{\textcolor[rgb]{0.00,0.23,0.31}{\textbf{#1}}}
\newcommand{\NormalTok}[1]{\textcolor[rgb]{0.00,0.23,0.31}{#1}}
\newcommand{\OperatorTok}[1]{\textcolor[rgb]{0.37,0.37,0.37}{#1}}
\newcommand{\OtherTok}[1]{\textcolor[rgb]{0.00,0.23,0.31}{#1}}
\newcommand{\PreprocessorTok}[1]{\textcolor[rgb]{0.68,0.00,0.00}{#1}}
\newcommand{\RegionMarkerTok}[1]{\textcolor[rgb]{0.00,0.23,0.31}{#1}}
\newcommand{\SpecialCharTok}[1]{\textcolor[rgb]{0.37,0.37,0.37}{#1}}
\newcommand{\SpecialStringTok}[1]{\textcolor[rgb]{0.13,0.47,0.30}{#1}}
\newcommand{\StringTok}[1]{\textcolor[rgb]{0.13,0.47,0.30}{#1}}
\newcommand{\VariableTok}[1]{\textcolor[rgb]{0.07,0.07,0.07}{#1}}
\newcommand{\VerbatimStringTok}[1]{\textcolor[rgb]{0.13,0.47,0.30}{#1}}
\newcommand{\WarningTok}[1]{\textcolor[rgb]{0.37,0.37,0.37}{\textit{#1}}}

\usepackage{longtable,booktabs,array}
\usepackage{calc} % for calculating minipage widths
\usepackage{caption}
% Make caption package work with longtable
\makeatletter
\def\fnum@table{\tablename~\thetable}
\makeatother
\usepackage{graphicx}
\makeatletter
\newsavebox\pandoc@box
\newcommand*\pandocbounded[1]{% scales image to fit in text height/width
  \sbox\pandoc@box{#1}%
  \Gscale@div\@tempa{\textheight}{\dimexpr\ht\pandoc@box+\dp\pandoc@box\relax}%
  \Gscale@div\@tempb{\linewidth}{\wd\pandoc@box}%
  \ifdim\@tempb\p@<\@tempa\p@\let\@tempa\@tempb\fi% select the smaller of both
  \ifdim\@tempa\p@<\p@\scalebox{\@tempa}{\usebox\pandoc@box}%
  \else\usebox{\pandoc@box}%
  \fi%
}
% Set default figure placement to htbp
\def\fps@figure{htbp}
\makeatother





\setlength{\emergencystretch}{3em} % prevent overfull lines

\providecommand{\tightlist}{%
  \setlength{\itemsep}{0pt}\setlength{\parskip}{0pt}}



 


\usepackage{booktabs}
\usepackage{caption}
\usepackage{longtable}
\usepackage{colortbl}
\usepackage{array}
\usepackage{anyfontsize}
\usepackage{multirow}
\definecolor{UBCblue}{RGB}{12, 35, 68}
\definecolor{forestgreen}{RGB}{34, 139, 34}
\definecolor{caliborange}{RGB}{192, 57, 43}
\definecolor{resulteal}{RGB}{20, 143, 119}
\definecolor{randomorange}{RGB}{230, 126, 34}
\definecolor{trueHRred}{RGB}{178, 34, 34}
\setbeamercolor{title}{fg=white,bg=UBCblue}
\setbeamercolor{frametitle}{fg=white,bg=UBCblue}
\setbeamercolor{structure}{fg=UBCblue}
\setbeamercolor{block title}{bg=UBCblue!85,fg=white}
\setbeamercolor{block body}{bg=UBCblue!5}
\setbeamercolor{block title alerted}{bg=caliborange!85,fg=white}
\setbeamercolor{block body alerted}{bg=caliborange!5}
\setbeamercolor{block title example}{bg=forestgreen!75,fg=white}
\setbeamercolor{block body example}{bg=forestgreen!5}
\usepackage{booktabs}
\setbeamertemplate{navigation symbols}{}
\makeatletter
\@ifpackageloaded{caption}{}{\usepackage{caption}}
\AtBeginDocument{%
\ifdefined\contentsname
  \renewcommand*\contentsname{Table of contents}
\else
  \newcommand\contentsname{Table of contents}
\fi
\ifdefined\listfigurename
  \renewcommand*\listfigurename{List of Figures}
\else
  \newcommand\listfigurename{List of Figures}
\fi
\ifdefined\listtablename
  \renewcommand*\listtablename{List of Tables}
\else
  \newcommand\listtablename{List of Tables}
\fi
\ifdefined\figurename
  \renewcommand*\figurename{Figure}
\else
  \newcommand\figurename{Figure}
\fi
\ifdefined\tablename
  \renewcommand*\tablename{Table}
\else
  \newcommand\tablename{Table}
\fi
}
\@ifpackageloaded{float}{}{\usepackage{float}}
\floatstyle{ruled}
\@ifundefined{c@chapter}{\newfloat{codelisting}{h}{lop}}{\newfloat{codelisting}{h}{lop}[chapter]}
\floatname{codelisting}{Listing}
\newcommand*\listoflistings{\listof{codelisting}{List of Listings}}
\makeatother
\makeatletter
\makeatother
\makeatletter
\@ifpackageloaded{caption}{}{\usepackage{caption}}
\@ifpackageloaded{subcaption}{}{\usepackage{subcaption}}
\makeatother

\usepackage{bookmark}
\IfFileExists{xurl.sty}{\usepackage{xurl}}{} % add URL line breaks if available
\urlstyle{same}
\hypersetup{
  pdftitle={Extreme Subgroup Effects Under a Uniform Treatment Benefit},
  pdfauthor={Larry F. León},
  hidelinks,
  pdfcreator={LaTeX via pandoc}}


\title{Extreme Subgroup Effects Under a Uniform Treatment Benefit}
\subtitle{A Simulation Study Based on Observed Trial Data}
\author{Larry F. León}
\date{2026-02-23}

\begin{document}
\frame{\titlepage}

\renewcommand*\contentsname{Table of contents}
\begin{frame}[allowframebreaks]
  \frametitle{Table of contents}
  \setcounter{tocdepth}{1}
  \tableofcontents
\end{frame}
\setcounter{tocdepth}{1}
\tableofcontents
}

\section{Motivation \& Overview}\label{motivation-overview}

\begin{frame}{The Core Question}
\phantomsection\label{the-core-question}
\begin{block}{Subgroup analysis in regulatory submissions}
When many subgroups are examined in a single trial, small subgroups can
produce hazard ratio estimates that look dramatic \emph{purely by chance}.
\end{block}

\begin{alertblock}{This study asks directly}
Under a DGM with a \textbf{perfectly uniform treatment benefit}
(true HR\,=\,0.70 for every patient), how often do standard Cox model
analyses produce extreme subgroup results?
\end{alertblock}

\begin{itemize}
\tightlist
\item
  Clinical subgroups of a given \(N\) are statistically
  indistinguishable from meaningless \emph{random} groupings of the same
  size
\item
  The \textcolor{randomorange}{\textbf{random}} benchmarks random
  combinations and should fluctuate around overall results with
  ``extremeness'' due to variability of small samples as opposed to true
  subgroup effects
\end{itemize}
\end{frame}

\begin{frame}[fragile]{Setting: Uniform Treatment Effect}
\phantomsection\label{setting-uniform-treatment-effect}
The DGM is built from a real trial (GBSG breast cancer, \(N = 686\))
using \texttt{generate\_aft\_dgm\_flex()}.

\begin{block}{Key feature --- constant log hazard ratio}
$$\psi^{0}(\mathbf{L}) = \log(\mathrm{HR}_{\mathrm{true}}) \quad \text{for all } \mathbf{L}.$$
\end{block}

\begin{itemize}
\tightlist
\item
  True overall HR = 0.70 (moderate benefit)
\item
  Every subgroup shares \emph{exactly} this HR
\item
  Any apparent subgroup finding is a false positive by construction
\end{itemize}

\begin{exampleblock}{Spline parameterisation}
\texttt{log.hrs = log(c(0.70, 0.70, 0.70))} --- uniform HR at biomarker
values $0$, knot $k$, and upper anchor $\zeta$.
\end{exampleblock}
\end{frame}

\section{Simulation Workflow}\label{simulation-workflow}

\begin{frame}{Pipeline Overview}
\phantomsection\label{pipeline-overview}
\small
\begin{center}
\begin{tabular}{clp{7cm}}
\toprule
\textbf{Step} & \textbf{Function} & \textbf{Purpose} \\
\midrule
1 & --- & Convert source data to consistent time scale \\
2 & \texttt{generate\_aft\_dgm\_flex()} & Fit Weibull AFT; build super population \\
3 & \texttt{check\_censoring\_dgm()} & Verify censoring matches source data \\
4 & \texttt{calibrate\_cens\_adjust()} & Find \texttt{cens\_adjust} via \texttt{uniroot} \\
5 & \texttt{simulate\_from\_dgm()} & Draw synthetic trials from calibrated DGM \\
6 & \texttt{coxph()} loop & Fit Cox in 56 subgroups; record HR and UB(HR) \\
7 & \texttt{summarize\_extreme\_sims()} & Compute all summary statistics \\
8 & \texttt{plot\_ub\_forest()} & Visualise ECI distributions \\
\bottomrule
\end{tabular}
\end{center}
\end{frame}

\begin{frame}{Workflow: DGM Setup \& Calibration}
\phantomsection\label{workflow-dgm-setup-calibration}
\begin{figure}[H]

{\centering 
\includegraphics[width=0.95\linewidth,height=0.85\textheight]{sim_flow_part1.pdf}

}

\caption{Workflow Part 1: DGM construction and censoring calibration.}

\end{figure}%
\end{frame}

\begin{frame}{Workflow: Simulation Loop \& Results}
\phantomsection\label{workflow-simulation-loop-results}
\begin{figure}[H]

{\centering 
\includegraphics[width=0.95\linewidth,height=0.85\textheight]{sim_flow_part2.pdf}

}

\caption{Workflow Part 2: Simulation loop and results.}

\end{figure}%
\end{frame}

\begin{frame}{Subgroup Definitions (1/3): ITT, Clinical, Continuous}
\phantomsection\label{subgroup-definitions-13-itt-clinical-continuous}
\begin{table}
\fontsize{8.0pt}{10.0pt}\selectfont
\begin{tabular*}{1\linewidth}{@{\extracolsep{\fill}}>{\raggedleft\arraybackslash}p{\dimexpr 22.50pt -2\tabcolsep-1.5\arrayrulewidth}>{\raggedright\arraybackslash}p{\dimexpr 120.00pt -2\tabcolsep-1.5\arrayrulewidth}>{\raggedright\arraybackslash}p{\dimexpr 120.00pt -2\tabcolsep-1.5\arrayrulewidth}}
\toprule
\# & Subgroup & Definition \\ 
\midrule\addlinespace[2.5pt]
\multicolumn{3}{>{\raggedright\arraybackslash}m{1\linewidth}}{\parbox{\linewidth}{ITT}} \\[2.5pt] 
\midrule\addlinespace[2.5pt]
1 & All Patients & All subjects (flag\_itt == 1) \\ 
\midrule\addlinespace[2.5pt]
\multicolumn{3}{>{\raggedright\arraybackslash}m{1\linewidth}}{\parbox{\linewidth}{Clinical}} \\[2.5pt] 
\midrule\addlinespace[2.5pt]
2 & Post-menopausal & meno == 1 \\ 
3 & Pre-menopausal & meno == 0 \\ 
4 & Grade 3 & grade == 3 \\ 
5 & Grade 1/2 & grade != 3 \\ 
\midrule\addlinespace[2.5pt]
\multicolumn{3}{>{\raggedright\arraybackslash}m{1\linewidth}}{\parbox{\linewidth}{Continuous}} \\[2.5pt] 
\midrule\addlinespace[2.5pt]
6 & Age (young) & age ≤ median \\ 
7 & Age (older) & age > median \\ 
8 & Age ≤ 50 & age ≤ 50 \\ 
9 & Age > 50 & age > 50 \\ 
10 & Tumour size (small) & size ≤ median \\ 
11 & Tumour size (large) & size > median \\ 
12 & Node-negative & nodes == 0 \\ 
13 & Nodes 1–3 & 0 < nodes ≤ 3 \\ 
14 & High nodes (>Q3) & nodes > Q3 \\ 
15 & PGR low & pgr ≤ median \\ 
16 & PGR high & pgr > median \\ 
17 & ER-low (<20) & er ≤ 20 \\ 
18 & ER-high (≥20) & er > 20 \\ 
\bottomrule
\end{tabular*}
\end{table}
\end{frame}

\begin{frame}{Subgroup Definitions (2/3): Two-Way Interactions}
\phantomsection\label{subgroup-definitions-23-two-way-interactions}
\small

\begin{table}
\fontsize{7.0pt}{8.0pt}\selectfont
\begin{tabular*}{1\linewidth}{@{\extracolsep{\fill}}>{\raggedleft\arraybackslash}p{\dimexpr 22.50pt -2\tabcolsep-1.5\arrayrulewidth}>{\raggedright\arraybackslash}p{\dimexpr 135.00pt -2\tabcolsep-1.5\arrayrulewidth}>{\raggedright\arraybackslash}p{\dimexpr 150.00pt -2\tabcolsep-1.5\arrayrulewidth}}
\toprule
\# & Subgroup & Definition \\ 
\midrule\addlinespace[2.5pt]
\multicolumn{3}{>{\raggedright\arraybackslash}m{1\linewidth}}{\parbox{\linewidth}{Meno × Grade}} \\[2.5pt] 
\midrule\addlinespace[2.5pt]
19 & Pre-meno / Grade 3 & meno==0 \& grade==3 \\ 
20 & Post-meno / Grade 3 & meno==1 \& grade==3 \\ 
21 & Pre-meno / Grade 1–2 & meno==0 \& grade!=3 \\ 
22 & Post-meno / Grade 1–2 & meno==1 \& grade!=3 \\ 
\midrule\addlinespace[2.5pt]
\multicolumn{3}{>{\raggedright\arraybackslash}m{1\linewidth}}{\parbox{\linewidth}{Meno × Age}} \\[2.5pt] 
\midrule\addlinespace[2.5pt]
23 & Pre-meno / Age ≤ 50 & meno==0 \& age≤50 \\ 
24 & Pre-meno / Age > 50 & meno==0 \& age>50 \\ 
25 & Post-meno / Age (young) & meno==1 \& age≤med \\ 
26 & Post-meno / Age (older) & meno==1 \& age>med \\ 
\midrule\addlinespace[2.5pt]
\multicolumn{3}{>{\raggedright\arraybackslash}m{1\linewidth}}{\parbox{\linewidth}{Meno × ER}} \\[2.5pt] 
\midrule\addlinespace[2.5pt]
27 & Pre-meno / ER-low & meno==0 \& er≤20 \\ 
28 & Pre-meno / ER-high & meno==0 \& er>20 \\ 
29 & Post-meno / ER-low & meno==1 \& er≤20 \\ 
30 & Post-meno / ER-high & meno==1 \& er>20 \\ 
\midrule\addlinespace[2.5pt]
\multicolumn{3}{>{\raggedright\arraybackslash}m{1\linewidth}}{\parbox{\linewidth}{Grade × Nodes}} \\[2.5pt] 
\midrule\addlinespace[2.5pt]
31 & Grade 3 / Node-neg & grade==3 \& nodes==0 \\ 
32 & Grade 3 / Node-pos & grade==3 \& nodes>0 \\ 
33 & Grade 1–2 / Node-neg & grade!=3 \& nodes==0 \\ 
34 & Grade 1–2 / Node-pos & grade!=3 \& nodes>0 \\ 
35 & Grade 3 / High nodes & grade==3 \& nodes>Q3 \\ 
\midrule\addlinespace[2.5pt]
\multicolumn{3}{>{\raggedright\arraybackslash}m{1\linewidth}}{\parbox{\linewidth}{Grade × ER}} \\[2.5pt] 
\midrule\addlinespace[2.5pt]
36 & Grade 3 / ER-low & grade==3 \& er≤20 \\ 
37 & Grade 3 / ER-high & grade==3 \& er>20 \\ 
38 & Grade 1–2 / ER-low & grade!=3 \& er≤20 \\ 
39 & Grade 1–2 / ER-high & grade!=3 \& er>20 \\ 
\midrule\addlinespace[2.5pt]
\multicolumn{3}{>{\raggedright\arraybackslash}m{1\linewidth}}{\parbox{\linewidth}{Grade × PGR}} \\[2.5pt] 
\midrule\addlinespace[2.5pt]
40 & Grade 3 / PGR low & grade==3 \& pgr≤med \\ 
41 & Grade 3 / PGR high & grade==3 \& pgr>med \\ 
42 & Grade 1–2 / PGR low & grade!=3 \& pgr≤med \\ 
43 & Grade 1–2 / PGR high & grade!=3 \& pgr>med \\ 
\midrule\addlinespace[2.5pt]
\multicolumn{3}{>{\raggedright\arraybackslash}m{1\linewidth}}{\parbox{\linewidth}{3-way}} \\[2.5pt] 
\midrule\addlinespace[2.5pt]
44 & Pre-meno / Yng / ER-low & meno==0 \& age≤50 \& er≤20 \\ 
45 & Pre-meno / Yng / ER-high & meno==0 \& age≤50 \& er>20 \\ 
46 & Post-meno / G3 / ER-low & meno==1 \& grade==3 \& er≤20 \\ 
47 & Post-meno / G3 / Node+ & meno==1 \& grade==3 \& nodes>0 \\ 
48 & G3 / Node+ / ER-low & grade==3 \& nodes>0 \& er≤20 \\ 
\midrule\addlinespace[2.5pt]
\multicolumn{3}{>{\raggedright\arraybackslash}m{1\linewidth}}{\parbox{\linewidth}{Size × Nodes}} \\[2.5pt] 
\midrule\addlinespace[2.5pt]
49 & Large / Node-pos & size>med \& nodes>0 \\ 
50 & Large / Node-neg & size>med \& nodes==0 \\ 
51 & Small / Node-pos & size≤med \& nodes>0 \\ 
52 & Small / Node-neg & size≤med \& nodes==0 \\ 
\bottomrule
\end{tabular*}
\end{table}
\end{frame}

\begin{frame}{Subgroup Definitions (3/3): Random Benchmarks}
\phantomsection\label{subgroup-definitions-33-random-benchmarks}
\begin{table}
\fontsize{9.0pt}{11.0pt}\selectfont
\begin{tabular*}{0.9\linewidth}{@{\extracolsep{\fill}}rlll}
\toprule
\# & Subgroup & Approx N & Construction \\ 
\midrule\addlinespace[2.5pt]
53 & random60 & \textasciitilde{}60 & Draw 60 random indices from trial \\ 
54 & random40 & \textasciitilde{}40 & First 40 of the random60 indices \\ 
55 & random20 & \textasciitilde{}20 & First 20 of the random60 indices \\ 
56 & random15 & \textasciitilde{}15 & First 15 of the random60 indices \\ 
\bottomrule
\end{tabular*}
\begin{minipage}{\linewidth}
\vspace{.05em}
Re-sampled every simulation. Nested: random15 ⊂ random20 ⊂ random40 ⊂ random60. No clinical meaning — true HR = 0.70 by construction.\\
\end{minipage}
\end{table}
\end{frame}

\section{Part I --- Building the Null
DGM}\label{part-i-building-the-null-dgm}

\begin{frame}[fragile]{Fitting \texttt{generate\_aft\_dgm\_flex()}}
\phantomsection\label{fitting-generate_aft_dgm_flex}
\begin{itemize}
\tightlist
\item
  Fits a \textbf{Weibull AFT model} to source data (GBSG)
\item
  Generates a large \emph{super population}
  (\(N_{\text{super}} = 5000\)) from which all future trials are drawn
\item
  \texttt{model\ =\ "null"} --- treatment effect identical for every
  patient
\item
  Censoring model estimated by AIC selection (Weibull / log-normal)
\end{itemize}

\begin{exampleblock}{R code}
\small\ttfamily
dgm\_null <- generate\_aft\_dgm\_flex(\\
\quad data = gbsg, model = "null",\\
\quad subgroup\_vars = NULL, ~~ n\_super = 5000)
\end{exampleblock}
\end{frame}

\section{Part II --- Censoring Diagnostics \&
Calibration}\label{part-ii-censoring-diagnostics-calibration}

\begin{frame}[fragile]{Censoring Consistency}
\phantomsection\label{censoring-consistency}
\begin{itemize}
\tightlist
\item
  Credible results require censoring in synthetic trials to match source
  data
\item
  \texttt{check\_censoring\_dgm()} compares three quantities:

  \begin{enumerate}
  \tightlist
  \item
    \textbf{Overall censoring rate} vs \texttt{rate\_tol} (default 10
    pp)
  \item
    \textbf{Censoring-time quantiles} among censored subjects
  \item
    \textbf{KM median censoring time} vs \texttt{median\_tol} (default
    25\%)
  \end{enumerate}
\end{itemize}
\end{frame}

\begin{frame}[fragile]{Calibrating \texttt{cens\_adjust}}
\phantomsection\label{calibrating-cens_adjust}
\texttt{calibrate\_cens\_adjust()} uses \texttt{uniroot} to find the
log-scale shift that equates a target statistic between reference and
simulated data.

\begin{block}{Objective function}
At each candidate \texttt{cens\_adjust}:
(1) generate synthetic subjects,
(2) extract target metric,
(3) return sim\_metric $-$ ref\_metric.
\end{block}

\begin{center}
\begin{tabular}{ll}
\toprule
Target & Equates \\
\midrule
\texttt{"rate"} & Overall censoring rate \\
\texttt{"km\_median"} & KM-based median censoring time \\
\bottomrule
\end{tabular}
\end{center}

\textbf{Monotonicity:} larger \texttt{cens\_adjust} \(\Rightarrow\)
longer censoring times.
\end{frame}

\section{Part III --- Generating Synthetic
Trials}\label{part-iii-generating-synthetic-trials}

\begin{frame}[fragile]{\texttt{simulate\_from\_dgm()}}
\phantomsection\label{simulate_from_dgm}
With the null DGM calibrated, each simulated trial is a single function
call:

\begin{enumerate}
\tightlist
\item
  Samples from super population with replacement
\item
  Assigns treatment (\(1{:}1\) randomisation)
\item
  Generates latent survival times from the AFT model
\item
  Applies calibrated censoring model
\item
  Truncates at \texttt{analysis\_time}
\end{enumerate}

Under the null DGM: \texttt{flag\_harm} \(\equiv\) 0 for every subject.
\end{frame}

\section{Part IV --- Simulation
Results}\label{part-iv-simulation-results}

\begin{frame}{Subgroups Under Study (56 total)}
\phantomsection\label{subgroups-under-study-56-total}
\footnotesize
\begin{columns}[T]
\begin{column}{0.48\textwidth}
\begin{tabular}{lp{3.5cm}}
\toprule
\textbf{Group} & \textbf{Examples} \\
\midrule
ITT & All Patients ($N\!\approx\!686$) \\
Clinical & Post-/Pre-menopausal, Grade \\
Continuous & Age, Size, Nodes, PGR, ER \\
Interaction & Meno $\times$ Grade, ER, Age; Grade $\times$ Nodes, ER, PGR \\
\bottomrule
\end{tabular}
\end{column}
\begin{column}{0.48\textwidth}
\begin{tabular}{lp{3.5cm}}
\toprule
\textbf{Group} & \textbf{Examples} \\
\midrule
3-way & Pre-meno/Young/ER-low, \ldots \\
Size $\times$ Nodes & Large/Node-pos, \ldots \\
\textcolor{randomorange}{\textbf{Random}} & \textcolor{randomorange}{random60, 40, 20, 15} \\
\bottomrule
\end{tabular}

\medskip

\begin{alertblock}{\footnotesize Key benchmark}
\textcolor{randomorange}{\texttt{random*}} subgroups have
\textbf{no clinical meaning}. Re-sampled each simulation.
\end{alertblock}
\end{column}
\end{columns}
\end{frame}

\begin{frame}{Summary Statistics}
\phantomsection\label{summary-statistics}
\begin{columns}[T]
\begin{column}{0.48\linewidth}
\begin{block}{HR point estimate distribution}
\begin{itemize}\small
\item Median + 1st--99th percentile (ECI)
\item Under truth, median $\approx 0.70$ for every subgroup
\item ECI width driven by subgroup $N$
\end{itemize}
\end{block}
\end{column}

\begin{column}{0.48\linewidth}
\begin{block}{Upper confidence bound UB(HR)}
\begin{itemize}\small
\item Upper end of the Cox 95\% CI
\item $\Pr(\text{UB} < 1.0)$: power to show benefit
\item Med UB $|$ UB$\ge$1: worst-case for the safety reviewer
\end{itemize}
\end{block}
\end{column}
\end{columns}
\end{frame}

\begin{frame}{Simulation: 1000 trials completed}
\phantomsection\label{simulation-1000-trials-completed}
\begin{itemize}
\tightlist
\item
  \textbf{1000} simulated trials from the null DGM (true HR = 0.70)
\item
  \textbf{56} subgroups evaluated per trial
\item
  Subgroup sizes range from \(N \approx 15\) (random15) to
  \(N \approx 686\) (ITT)
\end{itemize}

\begin{alertblock}{Re-use with new results}
\footnotesize
\texttt{res <- load\_sim\_results("new\_file.rds")}\\
\texttt{S <- summarize\_extreme\_sims(res\$sim\_hrs, res\$sim\_ubs, res\$sim\_ns, res\$subgroups)}\\
All downstream tables and plots update automatically.
\end{alertblock}
\end{frame}

\section{Results --- Forest Plots}\label{results-forest-plots}

\begin{frame}{HR Distributions: Single-Variable Subgroups}
\phantomsection\label{hr-distributions-single-variable-subgroups}
\begin{figure}[H]

{\centering \includegraphics[width=\linewidth,height=0.85\textheight,keepaspectratio]{extreme_subgroups_slides_files/figure-beamer/hr-forest-single-1.pdf}

}

\caption{HR distribution --- single-variable subgroups. True HR = 0.70
(red dashed).}

\end{figure}%
\end{frame}

\begin{frame}{HR Distributions: Combinations + Random Benchmarks}
\phantomsection\label{hr-distributions-combinations-random-benchmarks}
\begin{figure}[H]

{\centering \includegraphics[width=\linewidth,height=0.85\textheight,keepaspectratio]{extreme_subgroups_slides_files/figure-beamer/hr-forest-combo-1.pdf}

}

\caption{HR distribution --- combination and random-benchmark
subgroups.}

\end{figure}%
\end{frame}

\begin{frame}{UB(HR) Distributions: Single-Variable Subgroups
\{.shrink=5\}}
\phantomsection\label{ubhr-distributions-single-variable-subgroups-.shrink5}
\begin{figure}[H]

{\centering \includegraphics[width=\linewidth,height=0.82\textheight,keepaspectratio]{extreme_subgroups_slides_files/figure-beamer/ub-forest-single-1.pdf}

}

\caption{UB(HR) --- single-variable subgroups. Dotted: UB=1.0; Dashed:
true HR=0.70.}

\end{figure}%
\end{frame}

\begin{frame}{UB(HR) Distributions: Combinations + Random Benchmarks
\{.shrink=5\}}
\phantomsection\label{ubhr-distributions-combinations-random-benchmarks-.shrink5}
\begin{figure}[H]

{\centering \includegraphics[width=\linewidth,height=0.82\textheight,keepaspectratio]{extreme_subgroups_slides_files/figure-beamer/ub-forest-combo-1.pdf}

}

\caption{UB(HR) --- combination and random-benchmark subgroups.}

\end{figure}%
\end{frame}

\section{Key Findings}\label{key-findings}

\begin{frame}{HR Point Estimates}
\phantomsection\label{hr-point-estimates}
\begin{block}{Large subgroups ($N > 200$)}
Median HR tracks truth ($\approx 0.70$) closely; ECI is narrow.
\end{block}

\begin{alertblock}{Small subgroups ($N \lesssim 40$)}
\begin{itemize}
\item ECI for \texttt{random15}/\texttt{random20} spans $\sim\!0.25$ to $>1.0$
\item Clinical subgroups of similar size show \textbf{indistinguishable} ECI widths
\item Clinical rationale provides \textbf{no statistical protection}
\end{itemize}
\end{alertblock}

\begin{center}
\fcolorbox{UBCblue}{UBCblue!5}{\parbox{0.85\textwidth}{\centering
\textbf{Clinical subgroups of a given $N$ are statistically
indistinguishable from meaningless random groupings of the same size.}
}}
\end{center}
\end{frame}

\begin{frame}{Upper Bounds UB(HR)}
\phantomsection\label{upper-bounds-ubhr}
\begin{columns}[T]
\begin{column}{0.55\linewidth}
\begin{block}{The upper-bound distribution}
\begin{itemize}\small
\item \texttt{random40} ($N\!\approx\!40$): 1\% of trials produce UB(HR)\,$>$\,3.0--4.0
\item \texttt{random15} ($N\!\approx\!15$): values above 5.0 not rare
\item A reviewer seeing UB(HR)\,=\,3.5 in 15 patients cannot distinguish
      chance from a genuine safety signal
\end{itemize}
\end{block}
\end{column}

\begin{column}{0.42\linewidth}
\begin{exampleblock}{\footnotesize Illustrative 99th pctl of UB(HR)}
\centering\footnotesize
\begin{tabular}{lr}
\toprule
Subgroup & P99 UB \\
\midrule
ITT ($N\!\approx\!686$) & $<1.1$ \\
$N\!\approx\!200$ & $\sim\!1.5$ \\
$N\!\approx\!60$ & $\sim\!2.5$ \\
random40 & $\sim\!3.5$ \\
random20 & $\sim\!5.0$ \\
random15 & $>5.0$ \\
\bottomrule
\end{tabular}
\end{exampleblock}
\end{column}
\end{columns}
\end{frame}

\begin{frame}{Cox Model Works Correctly}
\phantomsection\label{cox-model-works-correctly}
\begin{block}{What the simulation confirms}
\begin{enumerate}
\item Cox model is \textbf{unbiased} --- median HR $\approx$ 0.70 for every subgroup
\item Confidence intervals have correct \textbf{coverage}
\item Variability driven \textbf{entirely by sample size}
\end{enumerate}
\end{block}

\begin{alertblock}{The interpretive hazard}
\begin{itemize}
\item UB(HR)\,=\,3.5 in a 15-patient subgroup is \emph{expected} under uniform benefit
\item Without simulation-based calibration, no way to quantify how ``extreme'' a result is
\item The \texttt{random*} benchmarks provide exactly this calibration
\end{itemize}
\end{alertblock}
\end{frame}

\section{Conclusions}\label{conclusions}

\begin{frame}[fragile]{Conclusions}
\phantomsection\label{conclusions-1}
\begin{enumerate}
\item
  \textbf{Small subgroups produce extreme results by chance.} Under
  uniform HR = 0.70, subgroups with \(N \lesssim 40\) regularly produce
  UB(HR) \(> 3.0\).
\item
  \textbf{Clinical rationale does not protect against noise.} Clinical
  subgroups have the \emph{same} HR/UB distribution as random benchmarks
  of that size.
\item
  \textbf{The Cox model is working correctly.} Wide CIs in small
  subgroups are accurate --- they reflect genuine uncertainty.
\item
  \textbf{Simulation-based calibration is essential.} Without a null-DGM
  benchmark, reviewers cannot quantify ``extreme'' for a given subgroup
  size.
\item
  \textbf{The \texttt{forestsearch} package} provides a complete,
  reproducible workflow.
\end{enumerate}
\end{frame}

\begin{frame}[plain]{}
\phantomsection\label{section}
\vfill
\begin{center}
{\Large\bfseries\color{UBCblue} Thank You}

\bigskip

{\large \texttt{forestsearch} R package}

\medskip

\url{https://larry-leon.github.io/forestsearch/}

\medskip

\texttt{pak::pak("larry-leon/forestsearch")}
\end{center}
\vfill
\end{frame}

\section*{Appendix}\label{appendix}
\addcontentsline{toc}{section}{Appendix}

\begin{frame}[fragile]{New-Results Workflow}
\phantomsection\label{new-results-workflow}
\footnotesize

\begin{Shaded}
\begin{Highlighting}[]
\FunctionTok{source}\NormalTok{(}\StringTok{"summarize\_extreme\_sims.R"}\NormalTok{)}

\CommentTok{\# ── Option A: Load existing results ─────────────────────────────}
\NormalTok{res }\OtherTok{\textless{}{-}} \FunctionTok{load\_sim\_results}\NormalTok{(}\StringTok{"results/extreme\_sims\_5000.rds"}\NormalTok{)}

\CommentTok{\# ── Option B: After running a new simulation loop ────────────────}
\FunctionTok{save\_sim\_results}\NormalTok{(sim\_hrs, sim\_ubs, sim\_ns, subgroups,}
                 \AttributeTok{file =} \StringTok{"results/extreme\_sims\_new.rds"}\NormalTok{,}
                 \AttributeTok{hr\_true =} \FloatTok{0.70}\NormalTok{,}
                 \AttributeTok{metadata =} \FunctionTok{list}\NormalTok{(}\AttributeTok{dgm\_seed =} \DecValTok{99}\NormalTok{, }\AttributeTok{n\_per\_trial =} \DecValTok{686}\NormalTok{))}

\CommentTok{\# ── Summarize (works with either option) ─────────────────────────}
\NormalTok{S }\OtherTok{\textless{}{-}} \FunctionTok{summarize\_extreme\_sims}\NormalTok{(res}\SpecialCharTok{$}\NormalTok{sim\_hrs, res}\SpecialCharTok{$}\NormalTok{sim\_ubs,}
\NormalTok{                             res}\SpecialCharTok{$}\NormalTok{sim\_ns, res}\SpecialCharTok{$}\NormalTok{subgroups)}

\CommentTok{\# ── All plots driven by S ────────────────────────────────────────}
\FunctionTok{plot\_hr\_forest}\NormalTok{(S, }\AttributeTok{panel =} \StringTok{"single"}\NormalTok{)}
\FunctionTok{plot\_hr\_forest}\NormalTok{(S, }\AttributeTok{panel =} \StringTok{"combo"}\NormalTok{)}
\FunctionTok{plot\_ub\_forest}\NormalTok{(S, }\AttributeTok{panel =} \StringTok{"single"}\NormalTok{)   }\CommentTok{\# \textless{}{-}{-} ub{-}forest{-}single}
\FunctionTok{plot\_ub\_forest}\NormalTok{(S, }\AttributeTok{panel =} \StringTok{"combo"}\NormalTok{)    }\CommentTok{\# \textless{}{-}{-} ub{-}forest{-}combo}
\end{Highlighting}
\end{Shaded}
\end{frame}

\begin{frame}{\texttt{summarize\_extreme\_sims()} API}
\phantomsection\label{summarize_extreme_sims-api}
\footnotesize

\begin{tabular}{lp{8.5cm}}
\toprule
\textbf{Field} & \textbf{Description} \\
\midrule
\texttt{hr\_q}, \texttt{ub\_q} & Quantile matrices [\texttt{n\_sg} $\times$ 3]: 1st, median, 99th \\
\texttt{mean\_n} & Mean subgroup $N$ across simulations \\
\texttt{pr\_ub\_lt1}, \texttt{pr\_ub\_ge1} & Power / failure rate per subgroup \\
\texttt{cond\_ub} & Conditional UB quantiles given UB $\ge$ 1 \\
\texttt{results\_tbl} & Formatted \texttt{data.frame} for printing \\
\texttt{ok}, \texttt{ok\_ub} & Logical: valid subgroups \\
\texttt{idx\_hr\_single}, \texttt{idx\_ub\_combo}, \ldots & Index vectors for forest panels \\
\texttt{cat\_ok}, \texttt{cat\_cols} & Category labels and colour map \\
\texttt{lab\_cmed\_ub} & Pre-formatted conditional median labels \\
\bottomrule
\end{tabular}
\end{frame}

\begin{frame}[fragile]{Comparing Two Simulation Runs}
\phantomsection\label{comparing-two-simulation-runs}
\footnotesize

\begin{Shaded}
\begin{Highlighting}[]
\NormalTok{S\_old }\OtherTok{\textless{}{-}} \FunctionTok{summarize\_extreme\_sims}\NormalTok{(...)   }\CommentTok{\# e.g. n\_sims = 1000}
\NormalTok{S\_new }\OtherTok{\textless{}{-}} \FunctionTok{summarize\_extreme\_sims}\NormalTok{(...)   }\CommentTok{\# e.g. n\_sims = 5000}

\NormalTok{comp }\OtherTok{\textless{}{-}} \FunctionTok{compare\_extreme\_sims}\NormalTok{(S\_old, S\_new,}
                              \AttributeTok{metric =} \StringTok{"both"}\NormalTok{,}
                              \AttributeTok{label1 =} \StringTok{"1000 sims"}\NormalTok{,}
                              \AttributeTok{label2 =} \StringTok{"5000 sims"}\NormalTok{)}
\FunctionTok{print}\NormalTok{(comp, }\AttributeTok{row.names =} \ConstantTok{FALSE}\NormalTok{)}
\CommentTok{\# Subgroup | N | MedHR | Delta\_MedHR | PrUBge1 | Delta\_PrUBge1 | MedUBcond}
\end{Highlighting}
\end{Shaded}

Useful for verifying convergence, comparing DGM variants, or assessing
sensitivity to censoring calibration.
\end{frame}

\begin{frame}[fragile]{Helper Utilities}
\phantomsection\label{helper-utilities}
\footnotesize

\begin{Shaded}
\begin{Highlighting}[]
\CommentTok{\# ── Optimal figure height ────────────────────────────────────────}
\NormalTok{fh }\OtherTok{\textless{}{-}} \FunctionTok{forest\_fig\_height}\NormalTok{(S, }\AttributeTok{panel =} \StringTok{"combo"}\NormalTok{, }\AttributeTok{type =} \StringTok{"ub"}\NormalTok{,}
                         \AttributeTok{row\_height =} \FloatTok{0.45}\NormalTok{, }\AttributeTok{overhead =} \FloatTok{1.5}\NormalTok{)}

\CommentTok{\# ── Save / load round{-}trip ───────────────────────────────────────}
\FunctionTok{save\_sim\_results}\NormalTok{(sim\_hrs, sim\_ubs, sim\_ns, subgroups,}
                 \AttributeTok{file =} \StringTok{"results/my\_run.rds"}\NormalTok{, }\AttributeTok{hr\_true =} \FloatTok{0.70}\NormalTok{,}
                 \AttributeTok{metadata =} \FunctionTok{list}\NormalTok{(}\AttributeTok{cens\_adjust =} \FloatTok{0.42}\NormalTok{))}

\NormalTok{res }\OtherTok{\textless{}{-}} \FunctionTok{load\_sim\_results}\NormalTok{(}\StringTok{"results/my\_run.rds"}\NormalTok{)}
\CommentTok{\# prints: Loaded: ... (5000 sims x 56 subgroups, saved ...)}

\CommentTok{\# ── Print method ──────────────────────────────────────────────────}
\FunctionTok{print}\NormalTok{(S)   }\CommentTok{\# shows n\_sims, n\_sg, hr\_true, then results\_tbl}
\end{Highlighting}
\end{Shaded}
\end{frame}




\end{document}
